
% !TeX root = ../../main.tex

This chapter examined two aspects of bringing social interaction into computational research on language acquisition: the simulation of child--caregiver dialogue and the learning effects of caregiver-derived feedback. The two studies addressed different questions, using different models and experimental procedures. Their joint contribution lies in identifying what each approach can establish, and what remains necessary for a more integrated account of learning through interaction.

\section{Simulation Fidelity and Learning Effects: Distinct Contributions}

The first study showed that language models can approximate some properties of child and caregiver language, but that this approximation depends on the level of analysis and the interaction setting. The few-shot prompting conditions improved several word- and utterance-level properties, especially in caregiver-like generation. However, discrepancies remained at the dialogue level: simulated caregivers showed higher semantic alignment and lower semantic diversity than the human reference. Extended interactions also increased several differences between model-generated and human language. These findings indicate that plausible responses to individual utterances do not, by themselves, establish that a model can reproduce the properties of sustained child--caregiver interaction.

The second study showed that caregiver-derived rewards have different effects on model production. Structural alignment produced the clearest improvements in generated grammaticality. Communicative feedback had smaller effects, with the more consistent benefits associated with clarification requests rather than acknowledgments. Affective feedback reduced grammaticality under both generation metrics. Semantic contingency reduced ChildCG scores at 1M and 10M, while its GEC scores were close to baseline at those scales and higher at 0.1M. Both rewards also changed other properties of generation. Importantly, the empirical rewards did not consistently improve performance on BLiMP or Zorro. The evidence therefore supports changes in productive behavior, without establishing broader grammatical generalization.

These contributions should be distinguished. The first study evaluates models as simulators of interaction; it does not investigate the learning process through which children acquire the simulated behaviors. The second uses models as controlled learners; it does not reproduce a complete conversation with a caregiver. In particular, its rewards are derived from natural conversations, not from the interlocutors evaluated in the first study. The benchmark therefore does not validate the learning environment of the reinforcement-learning study, and the reinforcement-learning results do not demonstrate improved performance on the interaction benchmark.

\section{Implications for Computational Models of Language Acquisition}

A first implication concerns the criteria used to evaluate computational models. Across both studies, improvement on one measure did not establish improvement on another. In the interaction benchmark, closer correspondence with human utterance-level properties did not eliminate differences in dialogue organization. In the learning study, increased grammaticality in free generation was not accompanied by consistent gains on controlled grammatical contrasts. Although these discrepancies concern different tasks, they raise a common methodological issue: the adequacy of a model must be assessed with respect to the particular behavior or learning claim under investigation. A single successful outcome is insufficient to establish a general account of language acquisition.

This is especially important when interpreting changes in model production. A model can produce a larger proportion of grammatical utterances by changing which utterance types it favors, without necessarily acquiring grammatical distinctions that transfer to other contexts. At the same time, minimal-pair benchmarks do not exhaust the grammatical regularities that a learner might acquire. As discussed in the second study, they may not directly target the errors most frequently encountered in the model's own productions. The disagreement between the two evaluations should therefore neither be ignored nor treated as conclusive evidence that no learning occurred. It identifies a more specific question: which regularities changed, and under what conditions do those changes generalize?

A second implication is that caregiver feedback should not be treated as a single, interchangeable source of support. Research on child language has described several ways in which caregivers respond to children's contributions \citep{clark_conversational_2020, nikolaus2023communicative}. The learning study makes this distinction experimentally explicit: optimizing different feedback signals leads to different production outcomes. The lower grammaticality observed under affective rewards, and under semantic rewards in part of the evaluation, does not establish that these forms of feedback are detrimental to children's development. Rather, it shows that the corresponding computational objectives do not uniformly favor grammatical production. Changes in length or lexical diversity motivate further questions about other dimensions of learning, but do not themselves demonstrate developmental benefits. In particular, the semantic and affective conditions retain more lexical entropy than the other reward families while remaining below the input-only baseline; the relevant comparison should therefore be stated explicitly.

The comparison also requires care with similarly named measures. Semantic alignment in the interaction benchmark and structural alignment in the learning study capture different properties and play different experimental roles. The former is used to assess correspondence with human dialogue, whereas the latter supplies a reward for model training. Their results cannot be combined into a general claim that more alignment is either beneficial or harmful. For acquisition research, the relevant question is which property of a caregiver response provides useful information, for which aspect of learning, and under which conditions.

\section{Limitations and Directions toward Integrated Interaction-and-Learning Models}

The studies leave several gaps between the present models and children's language acquisition. In the benchmark, age-conditioned generation reflects the behavior of already trained models asked to simulate children of different ages. It is not a developmental trajectory produced by accumulating child-like experience. In the learning study, controlled input sizes and feedback conditions make experimental comparisons possible, but do not establish that the learning procedure corresponds to children's learning. In addition, pretrained reward predictors bring information beyond the learner's caregiver-input corpus into the optimization process. The pretrained evaluators separately influence how outputs are scored; their linguistic knowledge enters the measurement procedure, rather than supplying an additional learning signal in this experiment. These distinctions limit how directly the results can be interpreted as accounts of human development.

The measures also provide incomplete descriptions of interaction and learning. The benchmark focuses mainly on linguistic properties and does not comprehensively assess whether a response appropriately addresses a child's intended meaning, communicative difficulty, or need for repair. Its few-shot prompts also introduce length constraints, so the effects of examples and instructions are not fully separated. Its multi-turn condition involves two imperfect simulators, so observed discrepancies may reflect their combined behavior rather than the limitations of either participant alone. Similarly, the learning study relies on automatic feedback annotations, learned reward predictions, and automatic grammaticality measures. Further validation should examine these components separately, including whether reward predictions remain reliable on utterances generated by the learner rather than drawn from the original corpus.

A further limitation concerns how feedback is represented. During reinforcement learning, the reward model is fixed and predicts a feedback value from a generated utterance. It does not produce a caregiver response that changes with the evolving conversational context or the learner's history. This abstraction allows particular signals to be isolated, but leaves out the linguistic content of an online response and the possibility of adjustments by both participants. Natural interaction provides a setting in which caregiver responses and children's subsequent contributions are related over time \citep{masek_where_2021, fusaroli_caregiver_2023}. Modeling that relation requires more than adding a fixed reward to otherwise independent productions.

An integrated approach could connect a learning model to a caregiver model whose responses are evaluated against natural interaction. The learner would update through repeated exchanges, while evaluation would track both the properties of the interaction and changes in the learner's language. Such a system would require separate tests of simulation fidelity and learning effects; improvements in one should not be assumed to validate the other. Matched controls could help isolate the contribution of contingent feedback, for example by disrupting the relation between an utterance and its feedback while preserving other properties of the learning setting.

Finally, stronger tests of acquisition would examine whether improvements persist beyond the utterances favored during training. This could involve identifying recurrent errors in child and model productions, constructing targeted contrasts, and testing transfer to new lexical items and combinations. Evaluation materials should be separated from those used to train reward models and optimize the learner. These tests would help distinguish changes in production preferences from the acquisition of regularities that can be applied more broadly. They would also provide a basis for investigating combinations of feedback signals, rather than assuming that the effects of isolated rewards will remain unchanged when they occur together.

\section{Chapter Conclusions}

This chapter provides two complementary contributions to computational research on language acquisition. The interaction benchmark establishes a framework for comparing simulated and natural child--caregiver dialogue, and identifies discrepancies that remain despite improvements in individual utterances. The feedback-learning study shows that signals derived from caregiver responses can selectively change model production, while leaving broader grammatical generalization unresolved.
